\chapter{Placement des ordres enfants}\label{ch:mx:child-order-placement}

Un algorithme doit acheter 800 actions dans la prochaine minute. Il peut se placer dans la file au cours acheteur et espérer, ou payer la fourchette tout de suite. Attendre fait gagner une demi-fourchette
s'il est exécuté, et coûte davantage que cela si le prix s'éloigne pendant l'attente. Ce chapitre formule ce choix comme un problème, mesure les politiques usuelles dans le marché simulé,
résout un modèle markovien des meilleurs cours par programmation dynamique, laisse un apprenant tabulaire retrouver la même réponse, et
cherche le déséquilibre des files au-delà duquel croiser immédiatement coûte moins cher.

\section{Le problème de placement}

\begin{definition}[Problème de placement des ordres]\label{def:mx:child-order-placement:problem}
Le \emph{problème de placement des ordres}\index{problème de placement des ordres} consiste à choisir, pour un \omterm{def:mx:the-almgren-chriss-framework:parent}{ordre enfant} qui doit être exécuté avant une échéance, où et quand
le poster (au meilleur cours, derrière celui-ci, à l'intérieur de la fourchette, ou en le croisant comme un \omterm{def:mx:the-limit-order-book:orders}{ordre au marché}) et comment le réviser, de manière à minimiser le
coût espéré par rapport au prix milieu au moment de la décision, pénalisé par sa dispersion.
\end{definition}

\begin{definition}[Risque de non-exécution, transaction de solde]\label{def:mx:child-order-placement:risk}
Le \emph{risque de non-exécution}\index{risque de non-exécution} est la probabilité qu'un ordre passif ne soit pas exécuté à son échéance, ainsi que le coût de ce qu'il faut alors faire. La
\emph{transaction de solde}\index{transaction de solde} est ce remède~: à l'échéance, la fraction non exécutée est annulée et envoyée comme un \omterm{def:mx:the-limit-order-book:orders}{ordre au marché},
au prix qu'il a atteint entre-temps.
\end{definition}

Un \omterm{def:mx:the-almgren-chriss-framework:parent}{ordre enfant} (chapitre~14) hérite d'une quantité et d'une échéance du calendrier du chapitre~16. Son coût relève de l'un de trois cas. Croiser immédiatement coûte
une demi-fourchette, plus ce que l'ordre consomme au-delà du meilleur cours. Se placer au cours acheteur et être exécuté rapporte une demi-fourchette~; sa probabilité d'exécution
relève de la file (la \omterm{def:mx:tick-size-queues-and-priority:queuevalue}{valeur de la position dans la file} du chapitre~6, la probabilité d'exécution passive du chapitre~18 du livre~7). Se placer sans être exécuté ajoute la
\omterm{def:mx:child-order-placement:risk}{transaction de solde} une fois le prix éloigné, c'est-à-dire la sélection adverse (livre~1, chapitre~1) des ordres qui ne s'exécutent pas. Harris et Hasbrouck (1996)
ont mesuré ce compromis sur des ordres SuperDOT du NYSE et constaté que les \omterm{def:mx:the-limit-order-book:orders}{ordres à cours limité} placés au cours affiché ou mieux surpassent les \omterm{def:mx:the-limit-order-book:orders}{ordres au marché}, même après une pénalité
pour les ordres qui ne se sont pas exécutés.

Le marché simulé de ce chapitre est \texttt{firm.agentmkt}, avec des fournisseurs de liquidité plus calmes que ceux du chapitre~16~: sur une séance de 30 minutes, la fourchette vaut
un tick dans 97.7~\% des cas, les meilleurs cours portent 14.6 lots (de 100 actions) en moyenne, le prix milieu bouge une fois toutes les 9.4 secondes et 137 actions se négocient par
seconde, de sorte que la tranche de 800 actions représente environ un dixième du volume d'une minute.

\section{Passif contre agressif}

Cinq politiques traitent les mêmes neuf tranches par séance (achats et ventes en alternance, à deux minutes d'intervalle), sur 24 séances avec le même flux d'ordres pour
chaque politique~:
\begin{itemize}
\item \emph{cross}~: un \omterm{def:mx:the-limit-order-book:orders}{ordre au marché} portant sur les 800 actions immédiatement~;
\item \emph{join}~: un \omterm{def:mx:the-limit-order-book:orders}{ordre à cours limité} au cours acheteur, laissé tel quel, soldé à l'échéance~;
\item \emph{behind}~: le même, un tick en dessous du cours acheteur~;
\item \emph{reprice}~: se placer au cours acheteur, et suivre celui-ci lorsqu'il s'éloigne~;
\item \emph{imbalance}~: croiser immédiatement si le déséquilibre des files du côté de l'ordre, $(B-A)/(B+A)$ avec $B$ la taille au cours acheteur et $A$ au cours vendeur
pour un achat (livre~7, chapitre~8), dépasse 0.5, une règle empirique courante~; sinon, reprice.
\end{itemize}

\begin{definition}[Règle de révision du prix]\label{def:mx:child-order-placement:reprice}
Une \emph{règle de révision du prix}\index{règle de révision du prix} indique quand un \omterm{def:mx:the-almgren-chriss-framework:parent}{ordre enfant} en attente est annulé et remplacé à un nouveau prix~: ici, dès que le
meilleur cours de son côté s'éloigne de lui, au prix de sa place dans la file.
\end{definition}

\begin{table}[ht]
\centering
\small
\begin{tabular}{@{}lrrrrr@{}}
\toprule
politique & coût (e.t.) & é.t. & face au croisement (e.t.) & exécutée passivement & soldée \\
\midrule
cross & 0.70 (0.04) & 0.53 & --- & 0~\% & 0~\% \\
join & $-0.24$ (0.06) & 0.89 & $-0.94$ (0.07) & 91.0~\% & 9.0~\% \\
behind & 0.76 (0.11) & 1.59 & 0.06 (0.11) & 16.7~\% & 83.3~\% \\
reprice & $-0.22$ (0.05) & 0.72 & $-0.92$ (0.06) & 97.3~\% & 2.7~\% \\
imbalance & $-0.19$ (0.05) & 0.78 & $-0.89$ (0.06) & 90.8~\% & 2.7~\% \\
\bottomrule
\end{tabular}
\caption{Exécution d'un \omterm{def:mx:the-almgren-chriss-framework:parent}{ordre enfant} de 800 actions dans un délai de 60 secondes dans le marché simulé~: 216 tranches, coût par action en ticks par rapport au prix milieu au moment de la décision (positif~: payé davantage),
son écart type, la différence appariée par rapport au croisement, et les actions exécutées passivement ou par la \omterm{def:mx:child-order-placement:risk}{transaction de solde}.
Données~: \texttt{mx\_place.sim\_study}.}\label{tab:mx:child-order-placement:policies}
\end{table}

Se placer au cours acheteur coûte 0.94 tick par action de moins que croiser, presque toute la fourchette~: la file exécute 91~\% de l'ordre dans la minute. Croiser est cher non pour sa demi-fourchette
mais pour les 3.2 lots du second niveau dans lequel il s'enfonce (0.70 tick en moyenne). Se placer un tick en dessous ne s'exécute presque jamais
(16.7~\%) et laisse 83~\% à la \omterm{def:mx:child-order-placement:risk}{transaction de solde} après que le prix a bougé~: c'est la pire politique et la plus dispersée. La révision du prix échange quelques
centièmes de tick de moyenne contre une dispersion plus faible (0.72 contre 0.89)~: l'ordre qu'elle sauve de la \omterm{def:mx:child-order-placement:risk}{transaction de solde} est celui dont le prix s'est éloigné.
La règle de déséquilibre a croisé au début de 6.5~\% des tranches et en a payé le prix~: 0.03 tick de plus que la révision du prix en moyenne.

\section{Gestion de la file et révision du prix}

Deux éléments déterminent un ordre en attente~: sa position dans la file et la question de savoir si la file sera encore là lorsque son tour viendra. Sa place est connue par le flux
(dans \texttt{firm.exchsim}, la fonction \texttt{working()} compte les actions devant lui dans la vue propre de l'agent sur le carnet)~; l'avenir de la file est ce que
prédit le déséquilibre~: une offre mince face à une demande profonde a tendance à être consommée en premier, et le prix remonte d'un tick lorsqu'on s'écarte d'un achat en attente. Lehalle et Mounjid (2017) ont construit
des stratégies à cours limité qui surveillent ce déséquilibre pour réduire la sélection adverse.

\begin{definition}[Placement tenant compte de la file]\label{def:mx:child-order-placement:qaware}
Le \emph{placement tenant compte de la file}\index{placement tenant compte de la file} conditionne chaque décision (attendre, réviser, croiser) à la position de l'ordre dans sa file et à la taille des files des deux côtés, et non
seulement aux prix et au temps restant.
\end{definition}

Une \omterm{def:mx:child-order-placement:reprice}{règle de révision du prix} a deux modes de défaillance. Trop empressée, elle abandonne sa position dans la file pour rien lorsque le cours vacille~; trop lente, elle laisse l'ordre derrière
le marché, là où \emph{behind} s'est arrêté. Les moteurs réels ajoutent un temps de repos minimal, une tolérance d'un tick avant de poursuivre, et une limite à la poursuite
avant de croiser~; sur ce marché, la règle simple capte déjà l'essentiel de la valeur.

\section{Placement comme problème de décision markovien}

Modélisons les meilleurs cours comme des files (les processus de naissance et de mort du livre~4, chapitre~8, et le modèle à files réactives de Huang, Lehalle et Rosenbaum,
dans lequel les taux dépendent de la taille des files). Un achat de $r$ lots se place derrière $n$ lots au cours acheteur~; le cours vendeur porte $a$ lots. Pendant un pas
$dt$, une vente au marché arrive avec la probabilité $\mu\,dt$ et prend un lot en tête de file (l'un des nôtres lorsque $n=0$, ce qui rapporte un demi-tick par rapport au prix milieu),
un lot devant nous s'annule avec la probabilité $\theta n\,dt$, le cours vendeur perd un lot avec la probabilité $(\mu+\theta a)\,dt$ et en gagne un avec $\lambda\,dt$.
Lorsque le cours vendeur se vide, le prix monte d'un tick et les lots restants recommencent dans des files neuves, à un tick de plus. À chaque pas, l'ordre peut
croiser~: $r$ lots contre $a$ au cours vendeur coûtent un demi-tick chacun pour les premiers $a$ et 1.5 pour le niveau suivant. La programmation dynamique (livre~4,
chapitre~9) donne la valeur $V_k(r,n,a)$ avec $k$ pas restants~:
\[
V_k=\min\bigl(C(r,a),\ \E[V_{k-1}\mid r,n,a]\bigr),\qquad V_0=C(r,a).
\]

\omcode{code/firm/placement/firm_placement.py}{147}{159}{Induction à rebours du problème de placement~: pour chaque nombre de lots restants, la valeur d'attendre un pas plutôt que de croiser maintenant, sur la grille des lots devant et des lots au cours vendeur.}\label{lst:mx:child-order-placement:dp}

Calibré sur le flux de la séance de calibration (\texttt{firm\_placement.calibrate}), $\lambda=0.99$ lots par seconde se placent sur chaque meilleur cours, chaque lot au
meilleur cours s'annule à $\theta=0.021$ par seconde, les \omterm{def:mx:the-limit-order-book:orders}{ordres au marché} prennent $\mu=0.69$ lots par seconde à chaque côté, et le second niveau porte 3.2 lots. Pour la
tranche de 800 actions, le plan ne croise jamais au départ~: sur le modèle, se placer vaut $-0.40$ tick par action contre 0.66 pour croiser. Croiser devient
le bon choix que pour de petits ordres pressés. Pour un lot unique, le plan croise lorsque le déséquilibre dépasse un seuil qui baisse à mesure que l'échéance approche
(\cref{fig:mx:child-order-placement:threshold})~: 0.86 avec quinze secondes restantes, 0.79 avec cinq, 0.66 avec trois. Pour huit lots, il ne croise jamais
au départ, avec quinze secondes ou soixante.

\begin{omfigure}
\begin{tikzpicture}
\begin{axis}[width=10cm,height=5.4cm,xlabel={\small secondes restantes},ylabel={\small seuil de déséquilibre},
  tick label style={font=\footnotesize},grid=major,grid style={gray!25},xmin=0,xmax=21,ymin=-0.6,ymax=1,table/col sep=comma]
\addplot[omDef,very thick,mark=*,mark size=1.4pt] table[x=seconds,y=cut]{figdata/microstructure/17-child-order-placement/threshold.csv};
\addplot[gray,dashed,domain=0:21] {0.5};
\node[gray,font=\scriptsize,anchor=south east] at (axis cs:20.5,0.5) {règle empirique};
\end{axis}
\end{tikzpicture}

\omcaption{Le déséquilibre des files au-delà duquel le plan croise un lot d'un coup, en fonction du temps restant~: la règle à seuil unique
qui se rapproche le mieux de la décision de programmation dynamique sur les files neuves du modèle. Données~: \texttt{mx\_place.threshold\_curve}.}\label{fig:mx:child-order-placement:threshold}
\end{omfigure}

\omcode{code/firm/placement/firm_placement.py}{330}{341}{Le plan comme politique de l'exécuteur de tranches~: les lots restants, les lots devant (ou un nouveau placement au fond de la file après un déplacement du prix) et les lots au cours éloigné, consultés dans la table de programmation dynamique.}\label{lst:mx:child-order-placement:policy}

Exécuté sur le marché simulé, les décisions du plan coïncident avec la révision du prix sur chaque tranche~: il n'a jamais jugé le croisement avantageux avant l'échéance.
Le modèle est pourtant trop optimiste sur l'attente~: il prévoit $-0.40$ tick par action là où le marché donne $-0.22$, car son prix ne bouge que
lorsqu'une file se vide, alors que les fondamentalistes simulés poussent le prix vers une valeur que le modèle ne voit pas.

Le seuil du modèle est-il réel~? Un lot avec cinq secondes restantes, 1\,824 tranches, croisement et révision appariés sur le même flux~:
la \cref{fig:mx:child-order-placement:urgent} montre l'écart selon le déséquilibre. Attendre coûte 0.26 tick de moins lorsque le déséquilibre est inférieur à $-0.5$~; l'économie
tombe à quelques centièmes au-dessus de zéro et n'est pas distinguable de zéro au-delà de 0.5 ($+0.03\pm0.07$ entre 0.5 et 0.75, $-0.11\pm0.08$
au-dessus). Le marché simulé confirme que l'avantage d'attendre disparaît là où le modèle le prévoit, mais à aucun déséquilibre croiser n'est significativement
moins cher.

\begin{omfigure}
\begin{tikzpicture}
\begin{axis}[width=10cm,height=5.4cm,xlabel={\small déséquilibre des files du côté de l'ordre au départ},
  ylabel={\small placer moins croiser (ticks)},tick label style={font=\footnotesize},grid=major,grid style={gray!25},
  xmin=-1,xmax=1,ymin=-0.5,ymax=0.3,table/col sep=comma]
\addplot[omDef,very thick,mark=*,error bars/.cd,y dir=both,y explicit] table[x=mid,y=diff,y error=se]{figdata/microstructure/17-child-order-placement/urgent.csv};
\addplot[black,dashed,domain=-1:1] {0};
\end{axis}
\end{tikzpicture}

\omcaption{Un lot avec cinq secondes restantes~: le coût de se placer au cours acheteur (avec révision du prix et \omterm{def:mx:child-order-placement:risk}{transaction de solde}) moins le coût de croiser d'un coup, par
déséquilibre des files au départ, avec erreurs types~; 1\,824 tranches appariées. Données~: \texttt{mx\_place.urgent\_study}.}\label{fig:mx:child-order-placement:urgent}
\end{omfigure}

Cont et Kukanov (2017) ont posé le problème de placement sur plusieurs plateformes, avec frais et rémunérations, comme un problème d'optimisation convexe portant sur le nombre d'actions à afficher
où et combien envoyer sous forme d'\omterm{def:mx:the-limit-order-book:orders}{ordres au marché}~; le chapitre~18 traite la partie consacrée au routage.

\section{Apprendre une politique de placement}

Le plan avait besoin des taux du modèle. Un apprenant peut trouver une politique à partir de la seule expérience~: le Q-learning tabulaire (le livre~12, chapitre~17, développe
l'apprentissage par renforcement) estime le coût futur $Q(s,a)$ de chaque état et de chaque action à partir d'épisodes simulés, puis agit de façon gloutonne. Nevmyvaka, Feng et
Kearns (2006) ont appliqué l'apprentissage par renforcement à l'exécution de transactions sur les données de carnet d'ordres du NASDAQ, avec un état formé du temps restant, de l'inventaire et de variables de marché.
Ici \texttt{firm\_placement.qlearn} exécute 200\,000 épisodes du carnet markovien avec l'état (lots restants, lots devant, lots au cours vendeur, tranche de cinq secondes
de temps) et les actions attendre ou croiser. Sa politique gloutonne coûte $-0.39$ tick par action sur 40\,000 épisodes neufs, contre $-0.40$ pour le plan~:
proche, dans un espace d'états dont il n'a visité que 4.0~\%. L'apprenant a retrouvé la réponse du plan (se placer et attendre) sans qu'on lui communique les taux~; il ne pouvait
pas mieux faire, et dans un espace d'états plus grand il lui faudrait bien plus d'épisodes, ou un approximateur de fonction, sujet du livre~12.

\section{Tutoriel~: exécuter un ordre enfant}\label{tut:mx:child-order-placement:tut}

\begin{tutorial}
\textbf{Objectif.} Comparer des politiques de placement pour une tranche d'une minute sur le marché simulé, résoudre le problème de placement par programmation dynamique sur un
carnet markovien calibré, et l'apprendre par Q-learning.
\textbf{État final~:} le \cref{tab:mx:child-order-placement:policies}, les \cref{fig:mx:child-order-placement:threshold,fig:mx:child-order-placement:urgent}
et les chiffres des sections 4 et 5.
\begin{tutsteps}
\item \textbf{Marché.} \texttt{mx\_place.MARKET}, \texttt{market\_stats()}~; \texttt{firm\_placement.calibrate(res)}.
\item \textbf{Politiques.} \texttt{SliceExecutor(slices, policy)} avec \texttt{Cross}, \texttt{Post(offset, reprice)}, \texttt{Imbalance},
\texttt{DPPolicy}~; \texttt{sim\_study()}.
\item \textbf{Plan.} \texttt{solve(model, r\_max, horizon\_s)}, \texttt{imbalance\_threshold(plan, r, k)}~; \texttt{dp\_thresholds()},
\texttt{threshold\_curve()}.
\item \textbf{Contrôle.} \texttt{urgent\_study()}~: un lot, cinq secondes, appariés par déséquilibre.
\item \textbf{Apprendre.} \texttt{qlearn(model, \dots)}, \texttt{q\_policy(q)}, \texttt{simulate(model, policy, \dots)}~; \texttt{model\_study()}~; tracer
avec \texttt{fig\_place.py}.
\end{tutsteps}
\textbf{Pour aller plus loin.} Donner plus de poids aux fondamentalistes (dérive de court terme plus forte) et trouver le déséquilibre à partir duquel croiser devient
moins cher sur le marché simulé~; ajouter une dérive du prix milieu au carnet markovien et observer le seuil du plan baisser~; remplacer la table Q par un petit réseau.
\end{tutorial}

\section{Construction~: placement}\label{bld:mx:child-order-placement:placement}

\begin{build}
\bfield{Rôle} La couche d'\omterm{def:mx:the-almgren-chriss-framework:parent}{ordres enfants} sous les calendriers du chapitre~16 et l'\omterm{def:mx:benchmark-algorithms:algo}{algorithme d'exécution} du chapitre~28~: politiques, modèle de carnet calibré, son plan optimal
et une référence apprise.
\bfield{Interface} \texttt{BookModel(lam, theta, mu, fresh, depth2)}, \texttt{calibrate(res)}, \texttt{cross\_cost(r, a, depth2)},
\texttt{solve(model, r\_max, horizon\_s, dt, n\_max, a\_max)} (\texttt{Plan.decide}, \texttt{.value}, \texttt{.fresh\_value}),
\texttt{imbalance\_threshold(plan, r, k)}, \texttt{simulate(model, policy, r, horizon\_s, \dots)}, \texttt{qlearn(\dots)}, \texttt{q\_policy(q)},
\texttt{SliceExecutor(slices, policy, check\_s)}, \texttt{Cross}, \texttt{Post}, \texttt{Imbalance}, \texttt{DPPolicy}.
\bfield{Règles} Lots de 100 actions~; coûts en ticks par rapport au prix milieu au moment de la décision, par action ou sommés sur les lots~; la \omterm{def:mx:child-order-placement:risk}{transaction de solde} à l'échéance fait
partie de toute politique passive~; les ventes en sont l'image miroir.
\bfield{Tests d'acceptation} \texttt{code/firm/placement/tests/}~: le croisement parcourt le carnet~; la valeur du plan égale son propre coût de Monte Carlo et ne fait pas moins bien que
croiser toujours ou ne jamais croiser~; la région de croisement grandit avec les lots devant et rétrécit avec le cours vendeur~; le Q-learning approche le plan~; l'exécuteur
mène à terme les tranches de croisement et de révision.
\bfield{Pour aller plus loin} Taux réactifs aux files~; un signal de dérive dans l'état~; plusieurs plateformes (chapitre~18).
\end{build}

\begin{omsources}
\item L.~Harris et J.~Hasbrouck, ``Market vs.\ limit orders: the SuperDOT evidence on order submission strategy'', \emph{Journal of Financial and
Quantitative Analysis} 31(2), 1996.
\item Y.~Nevmyvaka, Y.~Feng et M.~Kearns, ``Reinforcement learning for optimized trade execution'', \emph{Proceedings of the 23rd International
Conference on Machine Learning}, 2006.
\item W.~Huang, C.-A.~Lehalle et M.~Rosenbaum, ``Simulating and analyzing order book data: the queue-reactive model'', \emph{Journal of the American
Statistical Association} 110(509), 2015.
\item R.~Cont et A.~Kukanov, ``Optimal order placement in limit order markets'', \emph{Quantitative Finance} 17(1), 2017.
\item C.-A.~Lehalle et O.~Mounjid, ``Limit order strategic placement with adverse selection risk and the role of latency'', \emph{Market
Microstructure and Liquidity} 3(1), 2017.
\end{omsources}

\section{Exercices}

\begin{exercise}[$\star$]\label{exo:mx:child-order-placement:1}
Un achat de 8 lots croise un cours vendeur de 5 lots avec 3.2 lots au niveau suivant. Calculer son coût par action par rapport au prix milieu, avec une fourchette d'un tick.
\end{exercise}

\begin{exercise}[$\star$]\label{exo:mx:child-order-placement:2}
Avec les taux calibrés, calculer le temps moyen d'attente d'un lot au fond d'une file acheteuse de 15 lots avant que la file devant lui ne se vide, si la file devant
ne fait que diminuer (ventes au marché et annulations à $\theta$ par lot devant) et que le prix ne bouge pas~?
\end{exercise}

\begin{exercise}[$\star$]\label{exo:mx:child-order-placement:3}
Pourquoi reculer d'un tick derrière le cours acheteur coûte-t-il plus cher que croiser, alors que l'exécution ne coûte aucune fourchette lorsqu'elle a lieu~?
\end{exercise}

\begin{exercise}[$\star\star$]\label{exo:mx:child-order-placement:4}
Pourquoi le seuil du plan baisse-t-il à mesure que l'échéance approche~?
\end{exercise}

\begin{exercise}[$\star\star$]\label{exo:mx:child-order-placement:5}
Le modèle prévoit $-0.40$ tick par action pour se joindre à la file, et le marché simulé donne $-0.22$. Citer deux éléments que le modèle ne prend pas en compte et qui expliquent l'écart.
\end{exercise}

\begin{exercise}[$\star\star$]\label{exo:mx:child-order-placement:6}
Pourquoi la révision du prix réduit-elle davantage l'écart type du coût que sa moyenne~?
\end{exercise}

\begin{exercise}[$\star\star\star$]\label{exo:mx:child-order-placement:7}
\emph{Programmation.} Résoudre le plan pour un lot avec cinq secondes restantes. Dans quelle part des états de file neuve (pondérés par leur loi) le plan croise-t-il
immédiatement, avec les taux calibrés et avec $\lambda$ divisé par deux (moins de vendeurs qui rejoignent le cours vendeur, lequel se vide alors plus vite)~?
\end{exercise}

\begin{exercise}[$\star\star\star$]\label{exo:mx:child-order-placement:8}
\emph{Trouver la faille.} «~Nos exécutions passives coûtent moins une demi-tick par rapport au prix milieu au moment de l'exécution, donc l'exécution passive gagne une demi-tick sur chaque action.~»
\end{exercise}

\section{Problème~: Se joindre, reculer d'un cran ou croiser~?}

\begin{problem}[{Problème du week-end --- se joindre, reculer d'un cran ou croiser~?}]\label{pb:mx:child-order-placement:1}
Les \omterm{def:mx:the-almgren-chriss-framework:parent}{ordres enfants} d'un algorithme croisent à chaque fois la fourchette~; le desk demande s'ils devraient plutôt reposer dans le carnet et dans quel cas croiser est la bonne décision.

\medskip\noindent\textbf{Partie I --- Le problème.}
\begin{enumerate}
\item Énoncer le \omterm{def:mx:child-order-placement:problem}{problème de placement des ordres} et son coût.
\item Définir le \omterm{def:mx:child-order-placement:risk}{risque de non-exécution} et la \omterm{def:mx:child-order-placement:risk}{transaction de solde}.
\item Pourquoi reculer devant le cours acheteur est-il à peine possible sur le marché simulé~?
\item Décrire la fourchette, les files et le rythme du marché simulé.
\end{enumerate}

\medskip\noindent\textbf{Partie II --- Politiques sur le marché simulé.}
\begin{enumerate}[resume]
\item Donner le coût par action et son écart type pour croiser, se joindre à la file, reculer derrière le cours acheteur et réviser le prix.
\item Pourquoi croiser 800 actions coûte-t-il plus cher qu'une demi-tick~?
\item Quelle part de l'ordre se joint passivement à la file, et quelle part va à la \omterm{def:mx:child-order-placement:risk}{transaction de solde}~?
\item Combien a coûté la règle empirique sur le déséquilibre, et à quelle fréquence a-t-elle croisé~?
\end{enumerate}

\medskip\noindent\textbf{Partie III --- Le plan.}
\begin{enumerate}[resume]
\item Écrire le livre de Markov et ses transitions.
\item Écrire l'induction arrière.
\item Donner les taux calibrés.
\item Que fait le plan de la tranche de 800 actions, et quelle valeur attend-il~?
\item Donner le seuil de déséquilibre du plan pour un lot avec quinze, cinq et trois secondes restantes.
\item \emph{Énoncer le résultat nommé}~: le coût espéré par action de chaque politique de placement, et le seuil de déséquilibre des files au-delà duquel croiser
immédiatement coûte moins cher.
\item Que dit le marché simulé de ce seuil pour un lot avec cinq secondes restantes~?
\end{enumerate}

\medskip\noindent\textbf{Partie IV --- Apprentissage et réponse.}
\begin{enumerate}[resume]
\item Comment le Q-learning tabulaire estime-t-il la politique, et sur quel état~?
\item Quelle est sa proximité avec le plan, et quelle part de sa table a-t-il visitée~?
\item Pourquoi aurait-il des difficultés sur un carnet réel~?
\item Que diriez-vous au desk~?
\item En une phrase~: quand croiser immédiatement est-il la bonne décision~?
\end{enumerate}
\end{problem}

\section{Questions d'entretien}

\begin{interviewq}[$\star$ \iqroles{trader}]\label{iq:mx:child-order-placement:1}
Vous devez acheter 1\,000 actions dans la minute qui vient, sur une action dont la fourchette vaut une tick et dont les files sont profondes. \omterm{def:mx:the-limit-order-book:orders}{Ordre au marché} ou \omterm{def:mx:the-limit-order-book:orders}{ordre à cours limité}~?
\end{interviewq}

\begin{interviewq}[$\star\star$ \iqroles{chercheur}]\label{iq:mx:child-order-placement:2}
Comment le déséquilibre des files intervient-il dans la décision de croiser, et pourquoi~?
\end{interviewq}

\begin{interviewq}[$\star\star$ \iqroles{chercheur}]\label{iq:mx:child-order-placement:3}
Posez le placement d'\omterm{def:mx:the-almgren-chriss-framework:parent}{ordres enfants} comme un programme dynamique. Quel est l'état~?
\end{interviewq}

\begin{interviewq}[$\star\star$ \iqroles{développeur}]\label{iq:mx:child-order-placement:4}
Votre logique de révision du prix annule et remplace un ordre à chaque fois que le cours acheteur clignote. Qu'est-ce qui ne va pas, et que rajouteriez-vous~?
\end{interviewq}

\begin{interviewq}[$\star\star$ \iqroles{ingénieur ML}]\label{iq:mx:child-order-placement:5}
Utiliseriez-vous l'apprentissage par renforcement pour le placement d'ordres~? À quoi le compareriez-vous, et comment~?
\end{interviewq}

\begin{interviewq}[$\star\star\star$ \iqroles{trader}]\label{iq:mx:child-order-placement:6}
Votre taux d'exécution passive atteint 90~\% mais votre coût d'exécution par rapport à l'arrivée est pire que celui d'un collègue qui croise davantage. Comment est-ce possible~?
\end{interviewq}
